\documentclass[11pt]{article}
\usepackage[margin=0.85in]{geometry}
% Actual PDF check: keep source-credit filenames within the text column.
\setlength{\emergencystretch}{3em}
\usepackage{amsmath,amssymb,graphicx,enumitem}
\usepackage{xurl}
\usepackage[hidelinks]{hyperref}
\title{First to Four: Your Next Pick Can Lose the Series}
\author{An original documentary sports-data pilot proposal}
\date{27 September 2026}
\begin{document}
\maketitle

\begin{abstract}
One court, two balls, seven possible challenges. Win a challenge and you choose
what happens next; win four and the series is over. The proposed pilot turns that
simple transfer of choice into a readable, character-driven sports episode.
Basketball is the first version, with later variety adaptations possible. This
is a proposed campaign, not actual outreach, an agreement or a commissioned show.
Jesser, Ludwig and Atrioc are aspirational recipients or collaborators only:
there is no affiliation, endorsement, representation or contract.
\end{abstract}

\section{The format in one sentence}
\textbf{Every win earns a choice, but the choice does not guarantee the next win.}
Two players contest a best-of-seven challenge series, stopping at the first
four wins. The winner chooses the next challenge from Shooting, HORSE, 1v1,
Backboard, 3 Point, Rim and Bump. The actual order is decided during play; it is
not a producer's secretly fixed escalation. Establish the first picker and
whether repeats are allowed before filming because the owner rules leave those
details open. All observed historical starts are Shooting, not proof that they
must be. The title is a dramatic premise, not a causal research finding.

The editorial question after every challenge is concrete: what does the winner
choose now, and what does the other player think of that choice? Show the answer
before the next game, then let the result test the player's stated expectation.
The audience can track both the series score and the next-choice holder without
needing a rules lecture. Never insert an extra game after the genuine clinch.

\begin{figure}[htbp]
\centering
\includegraphics[width=\linewidth]{tex_images/sports/format_overview.png}
\caption{Original format diagram for a production briefing. Only a non-clinching winner picks another game; the first picker and repeats require agreement. The arrows show rules, not proven strategic advantage.}
\end{figure}

\section{Why these creator lanes are proposals, not promises}
\subsection{Basketball pilot: an aspirational Jesser fit}
A Jesser basketball pilot is the proposed first approach because the hook can
be shown with basketball action and a legible scoreboard, rather than explained
as an abstract format. The creative request would be a small original-footage
test, not permission to reuse an existing creator's footage, branding, voice or
audience. It does not assume availability, interest, compensation or a booking.
The script below uses an original documentary sports-data style; it does not
impersonate anyone's narration, catchphrases or production identity.

\subsection{Variety adaptations: aspirational Ludwig or Atrioc versions}
Ludwig or Atrioc are possible future variety-format recipients, not attached
talent. The transferable element is first-to-four with winner choice, not a
requirement that every version use basketball. A later adaptation might combine
agreed skill stations or game challenges with comparable time budgets, accessible
difficulty and visible scoring. It would need its own consented rules, rights
review and fresh data; basketball results do not predict another creator's
outcomes. Do not announce such a collaboration or imply their endorsement.

\section{A real log supplies stakes, not manufactured odds}
The supplied log covers \textbf{13 completed series} from 7 August through
16 September 2026. The read-only analysis engine confirms a \textbf{12--1}
series split for Liam and Owen and \textbf{50--23} challenge wins across
\textbf{73 played challenges}. It also verifies an \textbf{11-series} current
winning streak for Liam, equivalently Owen's 11-series losing streak within
this supplied history. Earlier unrecorded play is outside that claim.

The best listed 3 Point winner completion time is \textbf{1:22, or 82 seconds},
by Liam on 7 September, second listed session. It is not 85 seconds, a universal
record or a standardized head-to-head speed test. Some winner times are missing
or pending; losing times, attempts and rebounding assistance are not standardized.
Court and rim-height observations are absent for later sessions and must not be
filled by carrying August's observations forward.

\subsection{Two contrasting edits from the supplied chronology}
\begin{itemize}
\item \textbf{10 September, first listed session: OLLOLOL.} Owen wins the
opener; Liam takes the next two; Owen, Liam and Owen trade the next three;
Liam wins the seventh, Bump, to finish 4--3. This is a genuine seven-game
highlight, not a claim that every episode will reach a decider.
\item \textbf{16 September, first listed session: LLLOOL.} Liam goes up 3--0.
Owen takes 3 Point and HORSE to close to 3--2, but Liam wins Rim to finish 4--2.
Frame it as a comeback attempt stopped, not a completed comeback from 3--0.
It is later than the 10 September decider, so the decider is not the latest game.
\end{itemize}

Dates identify source entries. Same-day sessions follow supplied list order,
not known clock times. L and O encode the listed winners, not independently
logged choices. Bump and HORSE raw ``5-x'' strings have unconfirmed scoring
semantics; keep them as source annotations rather than animating invented points.

\subsection{What an honest data overlay can say}
Within-series adjacent-round records yield \textbf{60 transitions}: the previous
winner won the next round \textbf{35} times and lost it \textbf{25} times.
The descriptive hold fraction is $35/60$, about $58.3\%$. It is not a forecast
of the next pick's success, and there is no justified automatic 50-percent null
model when players may differ in skill. Actual picker decisions were not logged;
the previous winner is only the inferred picker under the house rule.

Challenge selection is not randomized, rounds stop at the clinch, and the sample
contains repeated play by the same two people. That combination creates
selection and clinch-truncation bias. A next-round loss does not prove that the
choice caused the loss; a hold does not prove picker advantage. Do not package
these fractions as betting odds, independent observations, causal effects,
statistical significance, a ranking of optimal picks or a validated prediction.
The live scoreboard should say ``observed in the supplied log'' and show the
denominator. With no justified predictive model, omit future win odds entirely.

\section{Ready-to-adapt email: not sent}
\textbf{Subject: Pilot proposal --- First to Four, where every win chooses the next challenge}

Hello Jesser team,

I am proposing a compact basketball pilot: two players, one court and two balls,
with a menu of seven challenges. The first player to four wins takes the series,
and each challenge winner chooses the next game. The question between games is
always visible: what do you pick when your opponent knows what you want?

There is already a small, owner-supplied log to explain the premise honestly:
13 completed series, including a 4--3 decider. It is source material for the
story, not a promise of competitive balance or a prediction for new participants.
The package includes a standalone house-rule article and original explanatory
diagrams, with unresolved mechanics marked for agreement before filming.

The proposed pilot is a 6--10 minute original-footage edit with a readable
scoreboard, actual on-camera choice announcements and a short post-game debrief.
Filming permissions, participant consent and safe ball-to-ball-only Bump rules
come first. Dinner and costume jokes are optional and never required.

Would your team be open to reviewing the format for a small test? Scope, rights,
budget, scheduling and any collaboration would require a separate agreement.
No endorsement or participation is assumed. Thank you for considering it.

\textit{Draft only. Add the sender's approved signature privately; do not publish
personal contact details or claim this message has been sent.}

\section{A roughly 90-second original opener}
Read naturally, then time a rehearsal; the intervals below are proposed edit
beats, not a guaranteed speaking rate. Show authentic footage from the fresh
session, or label historical reconstruction clearly. Never stage a result and
present it as a recorded competitive outcome.

\begin{enumerate}[start=1]
\item \textbf{0:00--0:20.} ``Winning this game gives you something dangerous:
the next choice. We have seven basketball challenges, two players, and one way
to end the day. First to four wins. After every challenge, the winner chooses
what we play next. You can pick your strength. You can try their weakness.
Neither guarantees the next point on this scoreboard.''
\item \textbf{0:20--0:45.} ``The old log gives us a reason to watch, not a
prediction. Thirteen series ended twelve to one. One of them still went all
seven games. Another reached three-nil, then three-two, before the comeback
stopped. Those are different stories. Today we are recording the choices,
the attempts, and any mercies instead of guessing what happened afterward.''
\item \textbf{0:45--1:10.} ``Here is the contract. We agree the first picker,
repeat policy and unresolved ties before we start. Bump is ball-to-ball only
on a clear court. Nobody is a target. We stop for bystanders or unsafe contact.
Dinner and costume jokes are opt-in, and anyone can say no. The scoreboard
tracks games won, not how much embarrassment somebody can tolerate.''
\item \textbf{1:10--1:30.} ``You will see the current challenge, the series score,
and who gets the next choice. When someone wins four, it ends. No bonus game
to rescue the edit. So before the first shot: what do you want to play, and
why do you think it helps you? Let us find out what the court says.''
\end{enumerate}

\section{Episode architecture: 6--10 minutes without a fake game seven}
\begin{figure}[htbp]
\centering
\includegraphics[width=\linewidth]{tex_images/sports/creator_episode.png}
\caption{Original proposed episode map. Timing windows are flexible: a sweep gets a shorter honest episode, while a real decider may sustain a longer one. This is a production plan, not footage already made.}
\end{figure}

\subsection{Broadcast-legible information}
Keep a persistent two-name scoreboard with four empty win slots per player,
the current challenge label and a next-choice indicator. Use text, shape and
color together so color is not the only cue. Make attempt counters specific:
five for Shooting, three for Backboard and four for Rim. In Rim, display the
BALL streak separately; three consecutive BALLs charge one miss. HORSE gets
letters, 1v1 gets one/two-point scores, and 3 Point gets five arc locations plus
half court. Bump gets clearly explained active roles, not an invented
first-to-five graphic based on an unconfirmed raw log string.

Show a compact diagram when a game first appears, then return to the action.
Use a brief freeze frame to explain an airborne catch/release or HIGH JUMP RULES
countback only when the relevant footage can support the explanation. Record
the preceding successful shot's attempt count; a countback cannot be truthfully
recreated from a reaction shot. Mark replays and keep the real score visible.

\subsection{A small, feasible production footprint}
Plan one permitted court, two suitable balls and a small crew: a locked wide
view for continuous action, a second operator for reactions/details, and a
scorekeeper who can also monitor boundaries when that does not compromise
safety. Crew size can change after a location assessment. Do not put a camera,
operator or unattended phone inside a rebound or landing lane. Capture clean
consented audio, slate challenge starts and announce choices before the attempt.
Allow breaks, water and a genuine option to stop; production pressure never
turns an unsafe jump or contact tactic into a requirement.

Obtain venue filming permission and participant releases appropriate to the
planned use. Keep bystanders out of frame or obtain permission; stop if a shared
court becomes occupied. No NBA footage rights, music rights, creator clips,
logos or broadcast rights are assumed. Use original footage and these original
diagrams, plus separately licensed assets if approved. A collaboration agreement
would need to settle ownership, approvals, compensation and distribution.

\section{The next session should improve the evidence}
The opportunity is not to add confident graphics to missing data. It is to
record a fresh session well enough that viewers can distinguish a measured
event from a participant's impression.
\begin{itemize}
\item \textbf{Choices and context:} actual picker, announced choice, stated
reason, first-picker method, repeats policy, court markers and measured rim
height. Keep missing measurements missing, not inferred from an earlier venue.
\item \textbf{Attempts and mercy:} per-location attempts, separate BALL events,
mercy offers, accepted mercies, agreed attempt allowances and countback handling.
Record opt-in social bits separately; do not infer a meal obligation from mercy.
\item \textbf{Timing:} start/stop definition, completion time, attempts,
assistance and interruptions. If comparing players, obtain both comparable
measurements; a winner-only time is not a head-to-head speed estimate.
\item \textbf{Privacy and publication:} explicit data opt-in separate from
filming permission, approved display names, limited recording access and an
agreed retention/deletion plan. Review identifying footage and free-text notes
before any release. Publish approved aggregate CSV from the fresh session,
not raw recordings, private contacts, passwords or financial information.
\end{itemize}

Document changes to ambiguous house rules before play. If someone withdraws
from a costume, interview or optional data release, do not edit that choice as
a sporting failure. Collect feedback on scoreboard comprehension, pacing and
the winner-choice hook after the pilot; propose success criteria in advance
rather than fabricating view counts or guaranteed conversion metrics.

\section{Sources, credits and the proposed next step}
\begin{enumerate}
\item \textbf{Rules source:} the separate \textit{First to Four: The Court-Series
House Rules} article, transcribing the complete owner-supplied mechanics with
transparent safety/consent overrides. Resolve its agreement checklist before
shooting. Its jokes are not compulsory production deliverables.
\item \textbf{Numerical source:} \nolinkurl{records/basketball/series.json}, analyzed
read-only by \nolinkurl{site_content/basketball.py}. The figures above are verified
against that supplied 13-series snapshot, not imported from a sports database.
They must be rechecked if the source changes. No inferred picker is relabeled
as an actual logged decision.
\item \textbf{Diagram credits:} both pitch figures are original, generated by
\nolinkurl{site_content/sports_diagrams.py}. The rules and pitch use PNG assets
with SVG companions; no TikZ, NBA footage or third-party creator art is required.
\end{enumerate}

The proposed next step is a private format review, a resolved pre-play agreement
and a consented fresh-session pilot. Only then consider approved outreach and
distribution. This article sends no email, performs no outreach, uploads no
recordings and auto-publishes no raw or private data. It describes a proposed
campaign, not an existing partnership or a completed production.

\end{document}